% Krieg--Ullrich locators below refer explicitly to the author manuscript;
% journal pagination is not inferred from the manuscript page numbers.
The certificate width is the established local diameter of information
\citep[Section~4.1]{novak2008}.
\citet[Theorems~2--3 and~5]{osipenko} also studies when adaptation cannot
improve noisy scalar recovery. Those criteria concern deterministic
reconstruction error and the existence of an optimal nonadaptive linear
reconstruction; the present criterion instead classifies randomized
penalties of the complete terminal information-fiber width.
Classical deterministic results already establish strong limits on
adaptivity: on convex symmetric classes with linear information,
\citet[Theorem~4.4]{novak2008} construct nonadaptive information with
no greater global diameter and at most twice the reconstruction radius.
For scalar targets the radius comparison is exact
\citep[Theorem~4.5]{novak2008}. A modern formulation of the factor-two
bound appears in \citet[Proposition~4.1]{krieg2025}.
Nonsymmetric convex uncertainty itself is also established:
\citet{novak1995} compares adaptive and nonadaptive optimal-recovery
errors using relations between Gelfand and Bernstein widths. The
classical objective is reconstruction error, whereas our randomized
penalty applies to the full terminal compatible interval. Our classification allows nonsymmetric convex sets and concerns randomized
penalties of the terminal local diameter, including a converse and exact
precision and moment separations. Its mean-risk interchange applies
\citet[Theorem~3.4]{sion1958}; the concave-penalty upper bound uses Jensen's
inequality.

For fixed deterministic scalar information, the optimal worst-case
reconstruction error is half the largest fiber width. The distinction here
concerns randomized local-width penalties and the order of averaging and
worst-case evaluation. Specifically, \citet[Section~2]{krieg2025}
and \citet[Section~9.1, Eq.~(9.15), author manuscript v2]{krieg2026}
study $\inf_A\sup_{y\in F}\mathbb E_\omega
\lVert T(y)-A_\omega(y)\rVert$, the error of a selected reconstruction.
Here the random transcript instead incurs the full compatible-interval
width $W_{Q_\omega(y)}(y)$.
Writing its scalar local radius as $r_{Q_\omega(y)}(y)=W_{Q_\omega(y)}(y)/2$,
our mean risk is $2\sup_y\mathbb E_\omega r_{Q_\omega(y)}(y)$.
For a randomized batch, this generally differs from
$2\mathbb E_\omega\sup_y r_{Q_\omega}(y)$, twice the expected global radius
discussed in \citet[Eq.~(2.4), Remark~6.1 and Section~6.1, author manuscript v2]{krieg2026}.
The midpoint reconstruction has pointwise error at most the local radius,
but its error at the actual record need not equal that radius. Consequently,
neither replacing actual reconstruction error by interval width nor
interchanging the expectation and supremum preserves the optimization
problem. The present theorem concerns the terminal local-width risk.
The fixed measurement dictionary is substantive: if every linear
functional were directly queryable at the same cost, a scalar affine
target could be determined by querying its linear part once.

